% Source presentations for register row N-c-001.

\subsection{Libkin's finite-model foundations}
\label{sec:p09-presentation-n-c-001-libkin}

Libkin fixes a relational vocabulary \(\sigma\) and writes
\(\operatorname{STRUCT}[\sigma]\) for its finite structures.  In this setting
the elementary compactness counterexample is already internal to finite
semantics.  For \(n\geq 1\), put
\[
 \lambda_n\ :=\
 \exists x_1\cdots\exists x_n
 \bigwedge_{1\leq i<j\leq n}x_i\neq x_j,
 \qquad
 T=\{\lambda_n:n\geq1\}.
\]
Every finite subset of \(T\) has a finite model: a structure with as many
elements as the largest requested index suffices.  No finite structure models
all of \(T\).  Thus compactness fails when both the premises and the conclusion
are interpreted only over finite structures
\cite[prop.~3.2]{Libkin2004}.  This argument does not challenge ordinary
first-order compactness, because an infinite structure models \(T\).

Libkin next develops finite Ehrenfeucht--Fra\"iss\'e games before turning to the
computability boundary.  Corollary~3.19 uses the finite game to show that graph
connectivity is not first-order definable, and Proposition~3.20 treats the
property of being a finite tree.  These are finite inexpressibility arguments:
for every quantifier rank, one supplies finite structures which the duplicator
cannot distinguish at that rank although the target graph property differs
\cite[pp.~37--40]{Libkin2004}.  They do not entail the stronger recursion-
theoretic conclusion below.

Definition~9.1 in the source distinguishes a sentence having some finite model
from a sentence true in every finite model.  With those meanings fixed,
Theorem~9.2 states the Trakhtenbrot boundary for every relational vocabulary
containing at least one binary relation symbol.  The proof encodes a finite
initial segment of a Turing-machine computation in a finite relational
structure, imposes the initial and local-transition conditions by a fixed
first-order theory, and proves the converse extraction from a satisfying
finite structure.  The displayed construction first uses several binary
relations and a constant.  Compression to a vocabulary with one binary
relation is assigned to Exercise~9.1 and is not part of the main-text proof
\cite[pp.~165--174]{Libkin2004}.  The resulting public theorem remains
Theorem~\ref{thm:trakhtenbrot}; no second theorem identity is introduced here.

Two consequences in the source sharpen the word ``undecidable.''  Corollary~9.3
says that finite validity is not recursively enumerable, so no effective
calculus can be both sound and complete for finite validity on the stated
vocabularies.  Corollary~9.4 rules out a recursive bound which, from a finitely
satisfiable sentence, bounds the size of one of its finite models
\cite[pp.~165--166]{Libkin2004}.  The latter is not merely a poor complexity
bound: no such recursive bound exists.

The same book records failure of finite Craig interpolation and finite Beth
definability in Exercises~3.12 and~3.13.  Those locations provide prompts and
hints rather than completed proofs \cite[pp.~41--42]{Libkin2004}.  Accordingly,
Theorem~\ref{thm:p09-finite-craig-beth} records the mathematical boundary, but
this source presentation does not misdescribe the two exercises as full
source-carried proofs.

\paragraph{Registered-edition boundary.}
The register also names the first edition of Ebbinghaus and Flum's
\emph{Finite Model Theory} (1995).  That edition remains a distinct authority:
the enlarged 2005 edition and Libkin's presentation cannot be used as its
edition-specific theorem text.  Because the 1995 theorem pages were not
available for verification, no notation, locator, proof, or Craig/Beth example
is attributed here to that edition merely from bibliographic metadata
\cite{EbbinghausFlum1995}.
